% Level 3, Week 9 -- How Do Machines Learn?
\lesson{9}{How Do Machines Learn?}{AI finds patterns in examples -- and copies their mistakes.}
\lessoninfo{Concepts: rules vs learning, features, training/test data, nearest neighbor, accuracy, bias, Markov text}{Tools: Thonny, turtle, random}{CSTA: 3A-IC-24, 3A-DA-12; AI4K12 Big Ideas 1--5}

\begin{objectives}
  \item The student can explain the difference between a program with hand-written rules and a program that learns from labeled examples. (L3.G7)
  \item The student can build a nearest-neighbor classifier and measure its accuracy on a separate test set. (L3.G7, L3.G2)
  \item The student can show how unbalanced or wrong training data leads to bad predictions, and why that is unfair. (L3.G7, L3.G9)
  \item The student can explain how a Markov text generator makes sentences, and why chatbots can be confidently wrong. (L3.G7, L3.G9)
\end{objectives}

\begin{materials}
  \item Print: workbook chapter 9, homework 9, and the \textbf{Project brief} (one per student, from \texttt{grade-4-6/3-advanced/project/}).
  \item On each Mac: in Thonny open \texttt{grade-4-6/3-advanced/exercises/week09/ex-monster-sorter/starter.py}.
  \item Unplugged: 11 index cards with quick monster doodles: 4 friendly (fluffy, few teeth) and 4 grumpy (bald, many teeth), labeled F or G; 3 unlabeled mystery cards.
  \item Prep (10 min): read this plan, make the cards, skim the Project brief.
\end{materials}

\begin{timeplan}
  \timeslot{0--6}{Warm-up}{Unplugged \enquote{Train the teacher}: you learn to sort monster cards from labeled examples.}
  \timeslot{6--14}{Concept}{Rules vs learning; features, labels, training and test data; the five big ideas of AI.}
  \timeslot{14--22}{Guided demo}{Type-along: Exercise 9.1 monster sorter (nearest neighbor on a chart).}
  \timeslot{22--42}{Hands-on build}{Pairs: Exercise 9.2 fair or biased (about 10 min), 9.3 Markov story machine (about 8 min). Bug hunt: bad training data. Swap at 10 min.}
  \timeslot{42--60}{Wrap-up}{AI ethics talk (4 min), exit ticket (1 min), then \textbf{final project kickoff} (about 13 min): brief, three options, milestones, pitch card.}
\end{timeplan}
\begin{teachernote}[If you are short on time]
Run 9.3's solution together instead of building it. Never cut the ethics talk or the project kickoff.
\end{teachernote}

\section{Warm-up (0--6 min)}
\begin{unplugged}[Train the teacher]
Say: \enquote{I am a robot who knows nothing about monsters. Teach me -- with examples only, no rules.} Students show you 8 labeled cards (\enquote{This one is friendly}). Hold up a mystery card and think aloud: \enquote{It looks most like \emph{this} friendly one, so I guess friendly.} Ask: \enquote{What did I look at?} (Fluffiness and teeth -- the \textbf{features}.) Then cheat on purpose: \enquote{What if you had only shown me grumpy monsters?} (Everything looks grumpy to the robot.)
\end{unplugged}

\section{Concept (6--14 min)}
\begin{teacherscript}
\enquote{Your programs so far follow \textbf{rules} you typed: \emph{if teeth > 5, say grumpy}. A \textbf{learning} program gets \textbf{examples} with the right answers (\textbf{labels}) and finds the pattern itself. What it measures are \textbf{features} -- here fluff and teeth. Our learner is the simplest: \textbf{nearest neighbor}. Put the examples on a chart; for a mystery monster, copy the label of the closest one. To check it, we hide some examples in a \textbf{test} pile, train on the rest and count how many test cards it gets right: its \textbf{accuracy}. Never test on cards it learned from -- that is seeing the exam answers first.}
\end{teacherscript}


\begin{teachernote}[The five big ideas of AI (AI4K12), for you]
Name them briefly (the workbook has the table). \textbf{1 Perception}: sensing with cameras and microphones. \textbf{2 Representation \& reasoning}: the world stored as data, like our (fluff, teeth). \textbf{3 Learning}: patterns from examples -- today's classifier. \textbf{4 Natural interaction}: talking with people is hard -- today's Markov generator. \textbf{5 Societal impact}: help or harm -- bias, privacy, deepfakes.
\end{teachernote}

\begin{analogy}[Learning birds with a field guide]
A child who has seen 100 robins and 2 crows calls every black bird \enquote{a dirty robin}. Not silly -- they learned from unbalanced examples. Machines can only be as fair as their examples.
\end{analogy}

\section{Guided demo (14--22 min)}
\begin{workedexample}[Exercise 9.1: monster sorter (type-along)]
\begin{enumerate}[itemsep=1pt]
  \item Read the \texttt{TRAINING} list aloud: each card is \texttt{(fluff, teeth, label)}, e.g. \py|(9, 1, "friendly")|.
  \item TODO 1: distance is the \textbf{steps distance}: the gap in fluff plus the gap in teeth, \py|abs(a[0] - b[0]) + abs(a[1] - b[1])|. Say: \enquote{Imagine a game piece on the grid from week 7. It can only move across or up and down. How many steps from (5, 2) to (6, 1)? One across, one down: 2 steps.} Count it on the chart with a finger. \py|abs()| means \enquote{the gap between two numbers}: \py|abs(2 - 5)| and \py|abs(5 - 2)| are both 3, so children do not need negative numbers. Check: from (0, 0) to (3, 4) is 3 + 4 = 7 steps. (The straight-line \enquote{ruler} distance needs Pythagoras, usually grade 8; it is an optional Extra in the README and gives the same guesses here.)
  \item TODO 2: \texttt{nearest} keeps the closest card so far -- \enquote{the same pattern as finding the smallest number in a list}.
  \item Run. Circles with F are friendly, squares with G grumpy; each white ? has a line to its nearest neighbor. Ask: \enquote{Do you agree with the guess for (5, 5), right in the middle?} Count together: three cards are exactly 4 steps away -- a tie! The program keeps the first one it found, (7, 3), so it says friendly.
\end{enumerate}
\end{workedexample}
\pyfile[firstline=30,lastline=60]{../exercises/week09/ex-monster-sorter/solution.py}

\section{Hands-on build (22--42 min)}
\begin{pairroles}[Data scientist and critic]
\textbf{Driver} types; \textbf{critic} asks \enquote{is that fair? could it be wrong?} at every result; swap at 10 minutes. 1 student: you are the critic. 3: add a \textbf{tester} who runs \texttt{check\_my\_work.py}; rotate. 4: two pairs who judge each other's Markov stories.
\end{pairroles}

\textbf{9.2} (10 min): \texttt{label\_counts}, \texttt{accuracy}; fair 100\% vs biased 50\%. \textbf{9.3} (8 min): \texttt{build\_chain}; read the stories.

\begin{bughunt}[Bad training data, bad predictions]
In 9.2 the code is correct but the biased predictions are wrong. Ask: \enquote{Is the bug in the code or in the data?} (The data: only one friendly example.) Then change the friendly training card \py|(7, 3, "friendly")| in \texttt{ALL\_CARDS} to \py|"grumpy"| and run again: the fair accuracy drops from 100\% to 75\% (3 of 4): the test card (6, 4) now copies the wrong label of its nearest card. One wrong label fools the guesses near it. (Change it back afterwards.) Strategy: \textbf{when a learning program is wrong, check the examples before the code.}
\end{bughunt}

\begin{freebuild}[Teach it something new (if time)]
Replace \texttt{TEXT} in 9.3 with your own poem (never someone's private messages!), or add mystery monsters in 9.1.
\end{freebuild}

\section{Wrap-up (42--60 min)}
\begin{safety}[AI ethics: wrong, private, fake and unfair]
\textbf{AI can be wrong:} our Markov machine writes sentences that \emph{sound} right but mean nothing; big chatbots can still invent facts confidently -- check important things with a trusted source or adult.
\textbf{Privacy with chatbots:} never type your real name, address, school, passwords, photos or secrets into a chatbot; what you type may be stored and used for training.
\textbf{Deepfakes:} AI can fake photos, voices and videos of real people -- ask \enquote{Who made this? Why?} before believing or sharing, and never make fakes of real people.
\textbf{Bias:} unbalanced examples gave unfair guesses today; systems that decide about people must be tested for fairness.
\end{safety}

\begin{formativecheck}[Exit ticket]
\textbf{Look for:} \enquote{it learns from examples}; blaming the data, not only the code, for bias.
\textbf{Question:} \enquote{A sorter says grumpy for every card. The test pile has 9 grumpy and 1 friendly. What is its accuracy? Is it a good sorter?}
\textbf{Answer:} 9 out of 10 (90\%), yet it never spots a friendly monster: high accuracy can hide unfairness to a smaller group. Accept \enquote{9 out of 10}.
\end{formativecheck}

\textbf{Final project kickoff (about 13 minutes)}
\begin{teacherscript}
\enquote{For three weeks you will \emph{Build Something Big}: a \textbf{grid game}, a \textbf{generative art app} or a \textbf{data or AI app} -- with classes, a smart data structure, a named algorithm, a saved file, tests and git. The brief has the full checklist.}
\end{teacherscript}
\textbf{(3 min)} Hand out the Project brief; read its checklist together. \textbf{(3 min)} Show the reference game \emph{Maze Muncher} (\texttt{grade-4-6/3-advanced/project/reference}; the Project teacher guide says how to run it) and point at each requirement: Player, Ghost and Chaser classes; a grid maze from a text file; Chaser uses breadth-first search with a queue; high scores saved. \textbf{(2 min)} Milestones on the board: \textbf{M1 Pitch} (this week), \textbf{M2 Design + core classes} (week 10), \textbf{M3 Features + 5 tests + playtest} and \textbf{M4 README + polish} (week 11), \textbf{Showcase} (week 12). \textbf{(1 min)} Point at the brief's \enquote{Pick your path} table: the \textbf{Core path} (a template plus one new class and an algorithm from class) is the recommended path for most students; the Full path is for extra challenge. \textbf{(4 min)} Students start the pitch card (workbook Exercise 9.4) and finish it as homework.

\begin{reflection}
What surprised you about how machines learn? Where did you meet AI this week? \emph{Then hand out homework 9 (fruit detective + pitch card).}
\end{reflection}

\section{Notes for the teacher}
\begin{commonmistakes}
  \item \textbf{Starting \texttt{best} at 0} instead of \py|training[0]|: then \texttt{best[2]} crashes. Start with a real card.
  \item \textbf{Markov chain overwrites}: \py|chain[word] = [next_word]| keeps only the last follower. Create the list once, then \texttt{append}.
\end{commonmistakes}

\begin{differentiation}
\supportpath Do 9.1 as a group, 9.2 in pairs; just run 9.3's solution. Act out \texttt{nearest} with the cards.
\extensionpath Key 9.3 on \emph{two} words (\py|(w1, w2)|) for more sensible stories.
\end{differentiation}

\begin{stretchchallenge}[Confusion table]
In 9.2, print a 2-by-2 table of true label vs guess. Which mistake does the biased sorter make most?
\end{stretchchallenge}
